\section{Структуры}

\begin{lstlisting}[caption={Структура \_IO\_FILE}, label={lst:IO_FILE}]
struct _IO_FILE
{
	int _flags;		/* High-order word is _IO_MAGIC; rest is flags. */

	/* The following pointers correspond to the C++ streambuf protocol. */
	char *_IO_read_ptr;	/* Current read pointer */
	char *_IO_read_end;	/* End of get area. */
	char *_IO_read_base;	/* Start of putback+get area. */
	char *_IO_write_base;	/* Start of put area. */
	char *_IO_write_ptr;	/* Current put pointer. */
	char *_IO_write_end;	/* End of put area. */
	char *_IO_buf_base;	/* Start of reserve area. */
	char *_IO_buf_end;	/* End of reserve area. */

	/* The following fields are used to support backing up and undo. */
	char *_IO_save_base; /* Pointer to start of non-current get area. */
	char *_IO_backup_base;  /* Pointer to first valid character of backup area */
	char *_IO_save_end; /* Pointer to end of non-current get area. */

	struct _IO_marker *_markers;

	struct _IO_FILE *_chain;

	int _fileno;
	int _flags2:24;
	/* Fallback buffer to use when malloc fails to allocate one.  */
	char _short_backupbuf[1];
	__off_t _old_offset; /* This used to be _offset but it's too small.  */

	/* 1+column number of pbase(); 0 is unknown. */
	unsigned short _cur_column;
	signed char _vtable_offset;
	char _shortbuf[1];

	_IO_lock_t *_lock;
	__off64_t _offset;
	/* Wide character stream stuff.  */
	struct _IO_codecvt *_codecvt;
	struct _IO_wide_data *_wide_data;
	struct _IO_FILE *_freeres_list;
	void *_freeres_buf;
	struct _IO_FILE **_prevchain;
	int _mode;
	int _unused3;
	__uint64_t _total_written;
	/* Make sure we don't get into trouble again.  */
	char _unused2[12 * sizeof (int) - 5 * sizeof (void *)];
};
\end{lstlisting}

\begin{lstlisting}[caption={Структура stat}, label={lst:stat}]
struct stat
{
	__dev_t st_dev;		/* Device.  */
	__ino_t st_ino;		/* File serial number.	*/
	__nlink_t st_nlink;		/* Link count.  */
	__mode_t st_mode;		/* File mode.  */
	__uid_t st_uid;		/* User ID of the file's owner.	*/
	__gid_t st_gid;		/* Group ID of the file's group.*/
	int __pad0;
	__dev_t st_rdev;		/* Device number, if device.  */
	__off_t st_size;			/* Size of file, in bytes.  */
	__blksize_t st_blksize;	/* Optimal block size for I/O.  */
	__blkcnt_t st_blocks;		/* Number 512-byte blocks allocated. */
	# ifdef __USE_XOPEN2K8
	/* Nanosecond resolution timestamps are stored in a format
	equivalent to 'struct timespec'.  This is the type used
	whenever possible but the Unix namespace rules do not allow the
	identifier 'timespec' to appear in the <sys/stat.h> header.
	Therefore we have to handle the use of this header in strictly
	standard-compliant sources special.  */
	struct timespec st_atim;		/* Time of last access.  */
	struct timespec st_mtim;		/* Time of last modification.  */
	struct timespec st_ctim;		/* Time of last status change.  */
	#  define st_atime st_atim.tv_sec	/* Backward compatibility.  */
	#  define st_mtime st_mtim.tv_sec
	#  define st_ctime st_ctim.tv_sec
	# else
	__time_t st_atime;			/* Time of last access.  */
	__syscall_ulong_t st_atimensec;	/* Nscecs of last access.  */
	__time_t st_mtime;			/* Time of last modification.  */
	__syscall_ulong_t st_mtimensec;	/* Nsecs of last modification.  */
	__time_t st_ctime;			/* Time of last status change.  */
	__syscall_ulong_t st_ctimensec;	/* Nsecs of last status change.  */
	# endif
	__syscall_slong_t __glibc_reserved[3];
};
\end{lstlisting}
